% from ega4-18-fr.tex
\oldpage[IV-4]{205}

La relation $\operatorname{H}_1(\mathbf f,M)=0$ implique donc $\operatorname{H}_0(f_n,\operatorname{H}_1(\mathbf f',M))=0$ et $\operatorname{H}_1(f_n,\operatorname{H}_0(\mathbf f',M))=0$. La première de ces deux relations signifie que l'on a $f_n\operatorname{H}_1(\mathbf f',M)=\operatorname{H}_1(\mathbf f',M)$ (III, 1.1.3.5). Or, par définition (III, 1.1.1), $\operatorname{H}_1(\mathbf f',M)$ est isomorphe à un quotient d'un sous-module $N$ de $M^{n-1}$; prenant sur $M^{n-1}$ la filtration des $M^j$ pour $j\leq n-1$, sur $N$ la filtration induite et sur $\operatorname{H}_1(\mathbf f',M)$ la filtration quotient de celle de $N$, on obtient sur $\operatorname{H}_1(\mathbf f',M)$ une filtration finie dont les quotients sont isomorphes à des quotients de sous-modules de $M$, donc sont séparés pour la topologie $\mathfrak J$-préadique par hypothèse. Comme la relation $f_n\operatorname{H}_1(\mathbf f',M)=\operatorname{H}_1(\mathbf f',M)$ implique \emph{a fortiori} $\mathfrak J\operatorname{H}_1(\mathbf f',M)=\operatorname{H}_1(\mathbf f',M)$, on déduit de l'hypothèse et des remarques précédentes que $\operatorname{H}_1(\mathbf f',M)=0$. L'hypothèse de récurrence prouve alors que la suite $\mathbf f'$ est M-régulière. D'autre part, la relation $\operatorname{H}_1(f_n,\operatorname{H}_0(\mathbf f',M))=0$ signifie que $f_n$ est $(M/(\sum_{i=1}^{n-1}f_iM))$-régulier, donc la suite $\mathbf f$ est M-régulière.

\begin{corollary}[19.5.2]
\label{IV.19.5.2}
Supposons que $A$ soit un anneau noethérien, que les $f_i$ ($1\leq i\leq n$) appartiennent au radical de $A$ et que $M$ soit un $A$-module de type fini. Alors les quatre conditions $a)$, $b)$, $b')$, $c)$ de (19.5.1) sont équivalentes.
\end{corollary}

On sait en effet alors ($0_{\mathrm I}$, 7.3.5) que tout $A$-module de type fini est séparé pour la topologie $\mathfrak J$-préadique.

\begin{proposition}[19.5.3]
\label{IV.19.5.3}
Soient
\[
  0\longrightarrow M'\longrightarrow M\longrightarrow M''\longrightarrow0
\]
une suite exacte de $A$-modules, $\mathbf f=(f_i)_{1\leq i\leq n}$ une suite M-régulière, $\mathfrak J=\sum_{i=1}^n f_iA$. Considérons les propriétés suivantes:
\begin{enumerate}
  \item[a)] La suite $\mathbf f$ est $M''$-régulière.
  \item[b)] La filtration $\mathfrak J$-préadique de $M'$ est induite sur $M'$ par la filtration $\mathfrak J$-préadique de $M$ (autrement dit, l'homomorphisme canonique $\operatorname{gr}_{\mathfrak J}^*(M')\to\operatorname{gr}_{\mathfrak J}^*(M)$ est injectif).
  \item[c)] L'homomorphisme canonique $M'/(\sum_{i=1}^n f_iM')\to M/(\sum_{i=1}^n f_iM)$ est injectif.
\end{enumerate}
Alors on a les implications $a)\Rightarrow b)\Rightarrow c)$, et $a)$ entraîne en outre que la suite $\mathbf f$ est $M'$-régulière. Si tout quotient d'un sous-module de $M''$ est séparé pour la topologie $\mathfrak J$-préadique, les conditions $a)$, $b)$ et $c)$ sont équivalentes.
\end{proposition}

Il est clair que $c)$ est conséquence de $b)$. Prouvons que sous les hypothèses de la dernière assertion, $c)$ entraîne $a)$: puisque $\mathbf f$ est M-régulière, on a, par (19.5.1), $\operatorname{H}_1(\mathbf f,M)=0$, d'où une suite exacte, portion de la suite exacte d'homologie
\[
  0\longrightarrow\operatorname{H}_1(\mathbf f,M'')\longrightarrow\operatorname{H}_0(\mathbf f,M')\longrightarrow\operatorname{H}_0(\mathbf f,M).
\]
Or, la condition $c)$ exprime que l'homomorphisme $\operatorname{H}_0(\mathbf f,M')\to\operatorname{H}_0(\mathbf f,M)$ est injectif (III, 1.1.3.5); elle équivaut donc à $\operatorname{H}_1(\mathbf f,M'')=0$, d'où, en vertu de l'hypothèse de séparation et de (19.5.1), le fait que $\mathbf f$ est $M''$-régulière.

Montrons ensuite que $a)$ entraîne $c)$ et le fait que $\mathbf f$ est $M'$-régulière, en procédant par récurrence sur $n$. Pour $n=1$, $f_1$ étant M-régulier est aussi $M'$-régulier et la propriété $c)$ n'est autre que le lemme (3.4.1.4). Pour $n>1$, l'hypothèse de récurrence montre que

